\section{Approach}
\label{sec:method}

Describe the system, algorithm, or study in enough detail for a technical
reader to understand the design decisions. Define notation before using it;
for example, a simple objective can be written as
\begin{equation}
  \hat{x} = \arg\min_x \; f(x) + \lambda g(x),
  \label{eq:objective}
\end{equation}
where $f$ is the primary objective and $g$ is a regularizer. Replace this
example with the equations and figures that are specific to your work.

Use a figure or table when it communicates structure more clearly than prose.
For figures stored as files, place them in the \texttt{figures/} directory and
include them with \texttt{\string\includegraphics}. Keep text inside figures
at least as large as the surrounding caption text.

\begin{table}[t]
  \centering
  \caption{A placeholder table. Replace it with a result from your work.}
  \label{tab:placeholder}
  \begin{tabular}{lcc}
    \toprule
    Configuration & Metric A & Metric B \\
    \midrule
    Baseline      & --       & --       \\
    Ours          & --       & --       \\
    \bottomrule
  \end{tabular}
\end{table}
